\documentclass[11pt,oneside]{report}
% ======================================================================
%  THE TRIADIC MEASUREMENT SUBSTRATE --- COMPLETE CORPUS (Volumes 1 & 2)
%  Aaron Switzer, 2026
%  Single-file build. Compile: pdflatex x3 (third pass settles the TOC).
% ======================================================================
% ======================================================================
%  TMS Volume 1 -- shared preamble
%  Compiles with pdfLaTeX both locally and on Overleaf (no exotic fonts).
% ======================================================================
\usepackage[T1]{fontenc}
\usepackage[utf8]{inputenc}
\usepackage{amsmath}
\usepackage{amssymb}
\usepackage{mathptmx}            % Times text + math (widely available)
\usepackage[scaled=0.90]{helvet} % sans for headings
\usepackage{courier}

\usepackage[letterpaper,margin=1in]{geometry}
\usepackage{amsthm}
\usepackage{mathtools}
\usepackage{bm}
\usepackage{microtype}
\usepackage{enumitem}
\usepackage{booktabs}
\usepackage{array}
\usepackage{longtable}
\usepackage{caption}
\usepackage{xcolor}
\usepackage{titlesec}
\usepackage{fancyhdr}
\usepackage[hidelinks]{hyperref}
\usepackage{tocbibind}

% ---- spacing -----------------------------------------------------------
\setlength{\parindent}{1.2em}
\setlength{\parskip}{0.35em}
\linespread{1.04}
\setlist{leftmargin=1.6em,topsep=0.25em,itemsep=0.15em,parsep=0pt}

% ---- theorem environments ---------------------------------------------
\newtheorem{theorem}{Theorem}[section]
\newtheorem{lemma}[theorem]{Lemma}
\newtheorem{corollary}[theorem]{Corollary}
\newtheorem{proposition}[theorem]{Proposition}

\theoremstyle{definition}
\newtheorem{definition}[theorem]{Definition}
\newtheorem{axiom}{Axiom}
\newtheorem{claim}[theorem]{Claim}
\newtheorem{principle}[theorem]{Principle}
\newtheorem{conjecture}[theorem]{Conjecture}
\newtheorem{prediction}[theorem]{Prediction}
\newtheorem{property}[theorem]{Property}
\newtheorem{postulate}[theorem]{Postulate}

\theoremstyle{remark}
\newtheorem{remark}[theorem]{Remark}
\newtheorem*{remark*}{Remark}

% ---- section heading look (unnumbered; author supplies the numbers) ----
\titleformat{\section}[hang]{\normalfont\large\bfseries\sffamily}{}{0pt}{}
\titleformat{\subsection}[hang]{\normalfont\normalsize\bfseries\sffamily}{}{0pt}{}
\titleformat{\subsubsection}[hang]{\normalfont\normalsize\bfseries\itshape}{}{0pt}{}
\titlespacing*{\section}{0pt}{1.4em}{0.5em}
\titlespacing*{\subsection}{0pt}{1.0em}{0.35em}
\titlespacing*{\subsubsection}{0pt}{0.8em}{0.3em}

% chapter (= each paper) heading
\titleformat{\chapter}[display]
  {\normalfont\sffamily\bfseries}
  {}{0pt}{\Huge}
\titlespacing*{\chapter}{0pt}{10pt}{24pt}

% ---- page-break quality: avoid stranded headings and orphaned/widowed lines ----
% (applies to both volumes via the shared preamble)
\widowpenalty=10000        % no single line at the top of a page
\clubpenalty=10000         % no single line at the bottom of a page
\displaywidowpenalty=10000 % same, around displayed equations
\brokenpenalty=4000        % discourage page breaks right after a hyphenated line
% titlesec: forbid a page break immediately after a heading (heading stranded at page foot)
\usepackage{needspace}
\usepackage{etoolbox}
\pretocmd{\section}{\Needspace*{4\baselineskip}}{}{}
\pretocmd{\subsection}{\Needspace*{3\baselineskip}}{}{}

% ---- running heads -----------------------------------------------------
\pagestyle{fancy}
\fancyhf{}
\fancyhead[L]{\small\itshape The Triadic Measurement Substrate}
\fancyhead[R]{\small\itshape\nouppercase{\rightmark}}
\fancyfoot[C]{\small\thepage}
\renewcommand{\headrulewidth}{0.4pt}
\fancypagestyle{plain}{\fancyhf{}\fancyfoot[C]{\small\thepage}\renewcommand{\headrulewidth}{0pt}}

% ---- abstract / keywords / references list ----------------------------
\newenvironment{paperabstract}
  {\begin{center}\begin{minipage}{0.88\textwidth}\small
   \noindent\textbf{Abstract.}\itshape\space}
  {\end{minipage}\end{center}\vspace{0.4em}}

\newcommand{\keywords}[1]{\noindent{\small\textbf{Keywords:} #1}\vspace{0.4em}}

% per-paper APA reference list (hanging indent), kept as authored
\newenvironment{reflist}
  {\par\small\setlength{\parindent}{0pt}%
   \setlength{\leftskip}{1.6em}\setlength{\parindent}{-1.6em}%
   \setlength{\parskip}{0.4em}%
   \raggedright\hyphenpenalty=50\exhyphenpenalty=50\relax
   \sloppy}
  {\par}

% convenience
\providecommand{\tightlist}{\setlength{\itemsep}{0pt}\setlength{\parskip}{0pt}}
\providecommand{\TMS}{\textsc{tms}}
\hyphenation{con-fir-ma-tion sub-strate}
\begin{document}

\thispagestyle{empty}
\begin{titlepage}
\centering
\vspace*{2cm}
{\sffamily\bfseries\Large The Triadic Measurement Substrate\par}
\vspace{0.6cm}
{\large $\bar{B} = 3\bar{\alpha}$\par}
\vspace{1cm}
{\sffamily\large Volume 2\par}
\vspace{1.4cm}
{\Large \bfseries Paper 9: The Total System: The Top of the Tower\par}
\vfill
{\large Aaron Switzer\par}
\vspace{0.4cm}
{\large 2026\par}
\vspace{1.2cm}
{\small\itshape Excerpted from the complete two-volume corpus, \emph{The Triadic Measurement Substrate}. Full corpus and reproducibility package: \url{https://doi.org/10.5281/zenodo.22035422}\par}
\end{titlepage}
\cleardoublepage
\pagenumbering{arabic}

% ---------------- Volume 2 / P9_TotalSystem ----------------
% ============================================================
% TMS Volume 2, Paper 9: The Total System
% Title line: The Top of the Tower
% Working draft. Front matter (Status/Notation/Preface) written LAST.
% Subsection-by-subsection sign-off rhythm. Current arc numbering.
% ============================================================

\chapter*{Paper 9}
\markboth{Paper 9}{Paper 9}
\addcontentsline{toc}{chapter}{Paper 9: The Total System: The Top of the Tower}

{\large\itshape The Total System: The Top of the Tower\par}\vspace{0.8em}

\noindent{\small\textbf{Keywords:} TMS, Total System, Promotion Operator, Codomain, Confirmation Tower, Final Coalgebra, Un-forming, Area Law, Re-forming, Generative Economy, Graded Subjecthood, Permanent Open Horizon}\par\vspace{0.6em}

\section*{Status of Results}
\addcontentsline{toc}{section}{Status of Results}

This paper uses the five-category labeling of Volume~1, extended through this volume, with each claim
labeled where it appears.

\textbf{Derived.} Results that follow from the five axioms, the Principle of Generative Economy, and
\ensuremath{U_{\mathrm{sh}}} by explicit construction or proof, or from results already derived in
Volume~1 and the earlier papers of this volume by explicit inheritance. That the promotion operator
\ensuremath{G(k)} has no codomain at the top level --- the level with no level above, where the
operation has an argument and no recipient --- is derived in this sense from the operator's
definition and the tower (\S\,II), and shown on the bit side to be the same fact as the top's shed
residue having no recipient field (\S\,IV, Appendix~B). That an un-forming reads out on two
orthogonal axes, the fate of the substrate and the fate of the boundary identity, partitioning into
four exhaustive cases, is derived from the boundary-reading of promotion and the orthogonality of
Paper~7 (\S\,III, Appendix~A). That the top re-forms --- carrying the compressible area-law surface
forward and shedding the volume interior to a manufactured ground --- is derived by inheritance from
the self-feeding ledger of Volume~1 and the Principle of Generative Economy, with no premise that
does not operate at an interior level (\S\,V, Appendix~B).

\textbf{Derived-conditional.} Results forced by Volume~1 and earlier-volume machinery, resting on
results internal to that machinery rather than on the bare axioms, and referee-checkable against
their sources. The per-fold shed of two bits and its locality, on which the area/volume split rests,
are inherited as Derived from Paper~3 and Paper~4 (\S\,IV, Appendix~B). The area-law exponent fixed
at two, the dimension of the now-surface agent manifold, is inherited from Volume~2 Paper~1 (\S\,IV,
Appendix~B).

\textbf{Structural correspondence.} Results in which a TMS construction reproduces the skeleton of a
known framework, never used as support for a TMS claim, only as a landing point after one. The top
closure's relation to a final coalgebra, and the reading of the equals-sign as isomorphism rather
than numerical identity (\S\,VII); the colimit hierarchy of Memory Evolutive Systems as the
categorical home of the tower (\S\,VII); and the conserved surface as an invariant of the kind
described in the identity-through-time literature (\S\,VII) are structural correspondences in this
sense.

\textbf{Predictions / Conjectures.} Falsifiable claims that follow from the framework but whose
closed forms or empirical tests are not established here. That an un-forming which is part pure
coupling and part binding-work partly survives and partly sheds, in a ratio set by how much of the
closure was relation rather than binding --- so that no un-forming is all-or-nothing except at the
two limits --- is a prediction in this sense (\S\,VI, Appendix~A); its quantitative test is not
established here.

\textbf{Permanent open horizon.} A question the framework is structurally unable to close, marked so
that no later claim assumes it closed. Whether the top closure is occupied --- whether there is
anything it is like to be the whole, in its standing or its re-forming --- is not settled by where
its promotion fails to go nor by what its surface carries forward, and is the sole permanent horizon
of this paper (\S\,V, Conclusion). This silence is load-bearing.

This five-way labeling applies throughout. Every claim in the body falls into exactly one category,
visible at the point of claim.

\section*{Notation}
\addcontentsline{toc}{section}{Notation}

Symbols are inherited from Volume~1 and the earlier papers of this volume. \ensuremath{G(k)} is the
promotion operator, carrying a dissolved level-\ensuremath{k} closure to a single symbol at level
\ensuremath{k{+}1} (Paper~8 \S\,IV); the \emph{tower} is the nested hierarchy of closures, each one a
symbol in the closure above it, with graded subjecthood at every level (Paper~2 \S\S\,VI--VII); the
\emph{top} is the level with no level above. \ensuremath{U_{\mathrm{sh}}} is the maximum-entropy
substrate field; \ensuremath{B(C)} is the constituent set of a closure \ensuremath{C} and
\ensuremath{\pi(C)} its boundary program, the symbol a level above reads. The per-fold shed is
\ensuremath{2} bits for binary-agent triads (Paper~3 \S\,IV). \ensuremath{I_{\mathrm{surface}}},
\ensuremath{I_{\mathrm{total}}}, and \ensuremath{I_{\mathrm{residue}}} are the area-scaling record, the
volume-scaling fold activity, and their difference (Appendix~B). An \emph{un-forming} is the
dissolution of a closure, read by the operator of Appendix~A.

\section*{Preface to Paper 9}
\addcontentsline{toc}{section}{Preface to Paper 9}

Paper 8 set a self under load and derived its dissolution as the conservation of promotion: the
record carried upward as resolution is lost. It closed by naming an exhausted region's
resolution-loss as the small-scale form of an event at the largest scale that it did not derive.
This paper derives that event. It asks what the promotion operator does at the one level the rest of
the volume does not reach --- the top of the tower, the level with no level above --- and finds the
operator there with an argument and no recipient. The result is a boundary condition: what becomes of
a closure that dissolves with nothing above to read it. The account is structural throughout, and is
kept apart, at the point each is made, from the question it does not address. Whether the whole is
occupied is not raised by the operator's behaviour at its boundary, and is held open where Paper~3
and Paper~8 left it.

\section*{I. The Top of the Tower}
\addcontentsline{toc}{section}{I. The Top of the Tower}

Paper 7 individuated a self by its closure (Paper~7 \S\,II); Paper 8 set that self
under load and derived its persistence condition and the event of its dissolution,
in which the promotion operator \ensuremath{G(k)} returns the dissolved identity as a
single symbol at the level above (Paper~8 \S\,IV, Conclusion). Paper 8 closed by
naming an exhausted region's resolution-loss --- its configurations conserved but no
longer readable as new, re-read under a new protocol to yield the next level --- as
the small-scale form of an event it located at the largest scale and did not derive
(Paper~8 \S\,VI). That event is the subject here.

Promotion is defined wherever a level above exists to receive the promoted symbol.
The confirmation tower of Paper 2 --- each closure one symbol in the closure above it
(Paper~2 \S\,VI), subjecthood graded across every level with none excluded
(Paper~2 \S\,VII) --- is the domain on which \ensuremath{G(k)} operates, and at every
interior level the operator has both an argument and a value: it reads a dissolved
closure and returns it one level up. The tower has two ends. Volume~1 fixed the lower
one, the triadic minimum below which no confirmation is defined (Volume~1, Paper~1
\S\,I). The upper end is the level with no level above, and there the value of
\ensuremath{G(k)} is not given by the interior account, because that account assumes a
recipient level which, at the top, does not exist. This paper determines what
\ensuremath{G(k)} does at that boundary, and derives the result from the operator and
the ledger already in hand rather than asserting it.

The determination is structural, and is kept distinct from a question it does not
address. Whether the top closure is occupied --- whether there is anything it is like
to be the whole --- is not settled by where its promotion goes, and is not made nearer
by it; it is the permanent open horizon of this paper, marked here so that no later
step assumes it closed (cf.\ Paper~3, Conclusion; Paper~8 \S\,II). What is asked is
narrower: given \ensuremath{G(k)} as Papers 7 and 8 define it and the area/volume
ledger of Volume~1 (Paper~5, Appendix~A), what becomes of a closure that dissolves
with no level above to read it.

\section*{II. Promotion Has No Recipient at the Top}
\addcontentsline{toc}{section}{II. Promotion Has No Recipient at the Top}

Promotion is a specific operation, not a general tendency. To promote a level-\ensuremath{k}
closure is for three level-\ensuremath{k} subsystems to meta-confirm a state at level
\ensuremath{k{+}1}, their outputs reaching the meta-interstice across the imprint depth
(Paper~8 \S\,III; Volume~1, Paper~4 \S\,5.3). The operation is defined by its target:
\ensuremath{G(k)} carries an argument at level \ensuremath{k} to a value at level
\ensuremath{k{+}1}, and the value exists only because level \ensuremath{k{+}1} exists to hold
it. At every interior level this is automatic --- the tower continues upward, and a dissolved
closure is read by whatever closure currently encloses it (Paper~8 \S\,IV, Conclusion).

The top level is defined as the one with no level above (Paper~2 \S\,VI). Substituting that
definition into the operation: \ensuremath{G(k)} at the top would carry its argument to level
\ensuremath{k{+}1}, and there is no level \ensuremath{k{+}1}. The operator does not return a
different value at the top; it has no value to return, because the codomain its definition
requires is empty. This is not a quantity that comes out small. It is the absence of the place a
quantity would be written. The top closure is therefore closed in the sense Paper 7 fixed --- it
is a region read against what it encloses --- but it is the unique closure that cannot be
promoted, not because promotion fails on it but because promotion is undefined where there is no
recipient level.

One route would restore a recipient and must be examined, because if it always succeeds there is
no true top and the result is vacuous. The route is the one Volume 1 used at region scale: an
exhausted system is re-read under a new protocol, and the re-read yields a level that did not
previously operate (Volume~1, Paper~5 \S\,IV; Paper~8 \S\,IV). If the top, on dissolving, were
always re-read into a pre-existing higher level, then that higher level was the top, and the
question recurs at it without terminating. The two cases are distinct and the distinction is
decidable from the definitions. A recipient level that already operates makes the closure beneath
it not the top; a closure with no operating level above it is the top by definition, and for it
the re-read cannot deliver an existing recipient, because there is none. What the re-read does
instead --- whether it manufactures a recipient rather than finding one --- is a separate
question, taken up in \S\,V; it is not a pre-existing codomain, and so it does not collapse the
present result. The top is the level at which \ensuremath{G(k)}, as defined, has no value:
promotion has an argument and no recipient. This is the structure recognized in the theory of coalgebras as a final object, situated below (\S\,``Where It Sits'').

\section*{III. Two Axes of an Un-forming}
\addcontentsline{toc}{section}{III. Two Axes of an Un-forming}

When a closure dissolves, two quantities are released, and Paper 7 already established that they
are independent. Promotion reads a closure by its boundary: the level above registers the programs
at the meta-tetrahedral corners, and the coarse-graining specifies those boundaries and not the
interior the closure encloses (Paper~7 \S\,II, Appendix~B; Volume~1, Paper~4 \S\,2.5.1). The
interior is populated by a separate operation, incorporation, and Paper 7 Appendix~B derives that
the two are functions of disjoint arguments --- the boundary of the surround alone, the interior of
the region alone --- so that no change in one is a change in the other (Paper~7 \S\,II,
Appendix~B.4). A dissolution therefore reads out on two axes that the orthogonality forbids
collapsing into one.

The first axis is the fate of the substrate the closure bound. By Axiom 3 the bits are conserved;
what becomes of them is whether they re-couple into a neighbouring closure or return to the unbound
field \ensuremath{U_{\mathrm{sh}}} (Paper~8 \S\,IV). The second axis is the fate of the boundary
identity --- the program the level above would read. By the promotion account, that identity
survives if and only if a level above reads it; this is what \ensuremath{G(k)} does when it returns
the dissolved closure as a single symbol one level up (Paper~8 \S\,IV, Conclusion). The two axes
are independent because, by Paper 7 Appendix~B.4, the interior fate (Axis~1, what the bound
substrate does) and the boundary fate (Axis~2, whether the identity is read) are functions of
disjoint arguments. The crossing of the two axes gives four cases, and each is realized:

A closure whose substrate re-couples and whose identity is re-read is grafted: it continues, read
by a new level above. A closure whose substrate sheds to \ensuremath{U_{\mathrm{sh}}} and whose
identity is promoted is the death of Paper 8 --- resolution lost, record carried upward (Paper~8,
Conclusion). A closure whose substrate re-couples but whose identity finds no reader is the case
Paper 8 did not treat: the bound matter persists and is read by some new closure, but as something
other than what it was, and the original identity, read by nothing at its level, lapses. The fourth
case, substrate shed and identity unread, is dissolution with no survivor on either axis.

That third case is the one the present paper needs, because it shows the two axes do not move
together. Its existence is established by a single instance: a closure whose every constituent
independently survives and re-couples, while the closure itself is read by nothing as what it was.
Paper 8's death does not cover it --- there the substrate sheds; here the substrate is conserved
entire and the identity lapses regardless. The independence is therefore not assumed but forced by
the orthogonality of Paper 7 Appendix~B.4: identity-fate is settled by whether a reader exists,
substrate-fate by where the bits go, and the two questions take disjoint inputs. The formal
statement of the un-forming operator and the proof that these four cases exhaust it are given in
Appendix~A. The hierarchy on which the operator acts has a categorical home in the colimit construction of Memory Evolutive Systems, situated below (\S\,``Where It Sits'').

\section*{IV. The Fold That Cannot Shed Upward}
\addcontentsline{toc}{section}{IV. The Fold That Cannot Shed Upward}

The promotion of \S\,II and III has a bit-side ledger, and reading it fixes what the top lacks. A
confirmation fold carries three agent-side loci to one bit-side outcome, and the map is not
injective: three binary loci span eight distinguishable inputs, one binary outcome two, so
distinction is shed of necessity, and the amount is fixed at two bits per fold (Paper~3 \S\,IV,
Appendix~B; Volume~1, Paper~5, Appendix~A). The fold's books therefore balance only when the shed
is counted: one confirmed bit retained, two bits shed to \ensuremath{U_{\mathrm{sh}}}, against
three loci's worth of distinction entering. This is the content of the master equation's
equals-sign --- not a numerical identity of one with three, which is false, but an informational
balance across a lossy map, exact once the two shed bits are entered (Paper~3 \S\,I, \S\,IV).

The shed is per-fold and local. Paper 3 fixes the shed at two bits for a single fold, and Paper 4
carries the same fixed quantity into the decay law, where each re-read sheds the same two bits, not
a growing or shrinking amount (Paper~4 \S\,IV). The balance is struck at each fold independently: a
fold's three loci, one retained bit, two shed bits close that fold's ledger and no other. There is
accordingly no quantity that accumulates across folds up the tower --- each level's fold sheds its
two bits to the level that holds \ensuremath{U_{\mathrm{sh}}} for it, and the books close there. The
decay law of Paper 4 is the statement that the shed recurs identically at each fold, never that a
deficit compounds (Paper~4 \S\,IV).

This locates the top's lack precisely. At every interior level the two shed bits are received by
the field \ensuremath{U_{\mathrm{sh}}} against which the next level operates; the shed has somewhere
to go, and the fold closes (Volume~1, Paper~5 \S\,IV). The top level is the one with no level above
(Paper~2 \S\,VI), and so the two bits its fold must shed have no level above to receive them. The
empty codomain of \S\,II is this same absence read on the bit side: \S\,II found that the promoted
symbol has no recipient level, and the ledger finds that the shed bits have no recipient field.
They are one fact. The top is not a fold that sheds a larger amount, nor one strained by
accumulation --- the shed is two bits as everywhere --- but a fold whose shed has no level above to
take it, exactly as its promotion has no level above to read it.

What the shed leaves behind, at the top as at every level, splits by the same exponent. The
accessible record of a region scales as a power of its linear size whose exponent Volume 1 derives
and fixes at two --- the same integer the per-fold shed carries, each forced from its own root
(Volume~1, Paper~5, Appendix~A; Paper~1 \S\,4.3; Paper~3 \S\,IV) --- while the total fold activity
scales as the volume. The difference, the volume in excess of that surface, is the residue that
does not survive as accessible record. The surface is the compressible, low-entropy part --- the
program identity a level above would read --- and the volume interior is the part shed to
\ensuremath{U_{\mathrm{sh}}}, uniform under its own measure. The computation of the split, and the
proof that the surface carries the program identity while the interior is the shed residue, is
given in Appendix~B.

\section*{V. The Top Re-Forms by the Ordinary Ledger}
\addcontentsline{toc}{section}{V. The Top Re-Forms by the Ordinary Ledger}

Section II left one route open: the shed residue might not require a pre-existing level above to
receive it, but instead constitute one. Volume 1 established this at region scale and did not mark
it as special to that scale. When a region exhausts, its shed residue is not discarded; the
centroids of the confirmation pattern at the new scale fall on the positions of the next-order
sublattice, so that each level's confirmed-bit positions are the next level's lattice positions
(Volume~1, Paper~5 \S\,IV). The residue does not await a ground to fall into --- it is the ground
the next level operates on. This is the operation at every interior rung, and \S\,IV fixed the
residue it acts on: the volume interior shed to \ensuremath{U_{\mathrm{sh}}}, uniform under its own
measure (Appendix~B).

At the top the same operation has the same inputs. The top's fold sheds its two bits and its volume
residue with no level above to receive them (\S\,IV), and the residue is, by the Volume 1 result,
the lattice on which a level not previously operating is confirmed (Volume~1, Paper~5 \S\,IV). No
premise enters here that did not operate at an interior rung: the residue-as-ground identity is
Volume 1's, the prohibition on a residue resting unread is the Principle of Generative Economy,
which forbids a configuration that could generate and does not, and which has been load-bearing
since Volume 1 (Volume~1, Paper~1 \S\,II; Paper~5 \S\,II). The top does not promote into a level
above, because there is none; it does not rest, because the Principle forbids it; it re-forms,
because the residue it sheds is the ground a next level is confirmed upon, by the operation that
runs at every rung. The re-forming is therefore not a new postulate at the top but the interior
ledger applied where its recipient is absent --- the residue manufacturing the level its shedding
had no level above to supply.

What re-forms carries forward what promotion always carries and sheds what promotion always sheds.
By \S\,IV and Appendix~B the surface is the compressible record --- the program identity --- and the
volume interior is the shed residue. A re-forming reads the surface as the identity carried and
returns the interior to \ensuremath{U_{\mathrm{sh}}} as the ground the next level operates on. The
top is no exception: its surface --- the structure of the whole as a closed pattern --- is the part
carried, and its interior is the part shed to the new ground. The re-forming conserves the boundary
and not the interior, which is what promotion does at every level (Paper~7 \S\,II, Appendix~B). The
whole re-forms with its surface carried and its interior returned to the field, by the same lossy
operation that promotes a self at its death. The conserved surface is an invariant of the kind described in the identity-through-time literature, and the question of whether the whole is a subject has been taken up directly in the cosmopsychist program; both are situated below (\S\,``Where It Sits'').

This is a structural result and nothing more is claimed for it. Whether the re-formed whole, or the
whole that preceded it, is occupied --- whether there is anything it is like to be either --- is not
settled by the conservation of a surface across the re-forming, and is not made nearer by it. The
ledger balances without that term: occupancy appears in neither the shed, nor the residue, nor the
surface carried across. The horizon stands where Paper 3 and Paper 8 left it (Paper~3, Conclusion;
Paper~8 \S\,II), and the silence is load-bearing.

\section*{VI. Predictions and Directions}
\addcontentsline{toc}{section}{VI. Predictions and Directions}

What this paper leaves open, it leaves open precisely. The no-codomain result, the two axes of an
un-forming, and the re-forming by the interior ledger close here. Three things do not, and are
marked.

The two axes of \S\,III predict a spectrum. A closure that is part pure coupling and part
binding-work partly survives and partly sheds: the independently-closing constituents persist and
re-couple, the bound-only structure sheds to \ensuremath{U_{\mathrm{sh}}}, in a ratio set by how
much of the closure was relation and how much was binding. This is a structural consequence of the
un-forming operator (\S\,III, Appendix~A), falsifiable in the form of its prediction --- that no
un-forming is a clean all-or-nothing between survival and shedding except at the two limits --- but
its quantitative test against a modelled closure is not established here, and awaits the same
simulation pinning Paper 8 reserved for the stressor constants (Paper~8 \S\,VI).

The residue magnitudes are inherited as open. The form of the area/volume split is fixed (\S\,IV,
Appendix~B), but the ratio of the conserved surface to the shed interior, at cosmic scale, depends
on the same two rates Paper 4 carried as underived in the cosmic-residue balance (Paper~4 \S\,IV).
The structural result --- that a surface is carried and an interior shed --- does not depend on
their magnitudes; the magnitudes are not new freedoms introduced here.

One result points forward. The conserved surface of \S\,V --- the program identity a re-forming
carries --- is the structure on which Paper 10 will rest a question this paper does not raise:
whether what is conserved across an un-forming has standing. Paper 10 develops that toward moral
standing, importing its bridge premises by hand and deriving none; here the surface is established
only as what the ledger conserves, with no claim that conservation confers standing. Where the
structural account ends and the imported premise begins is marked at the seam, in Paper 10, and
nothing of its conclusion is settled here.

\section*{VII. Where It Sits}
\addcontentsline{toc}{section}{VII. Where It Sits}

The no-codomain structure of \S\,II is recognized in the theory of coalgebras. A final coalgebra
is the object to which every coalgebra of its functor maps by a unique morphism, and by Lambek's
theorem its structure map is an isomorphism --- the object is isomorphic to its own image under the
functor (Lambek 1968; Rutten 2000). The top closure stands in that relation: every closure maps
toward it, and it has no level above because it is isomorphic to its own image, its own above. The
correspondence also fixes the reading of the equals-sign of \S\,IV --- the relation is isomorphism,
not numerical identity, the distinction the categorical setting draws between a fixed point and an
equality. The result is derived here from the operator and the ledger; the coalgebraic structure is
where it is recognized, not where it comes from.

The hierarchy the operator acts on has, in turn, a categorical home in Memory Evolutive Systems,
which model a multi-level system as a category in which each higher-level object is the colimit of a
pattern of linked lower-level objects, the higher level acquiring an individuation independent of
its constituents (Ehresmann and Vanbremeersch 2007). The colimit is the binding-into-one called
promotion here, and the independent individuation is the re-routability of \S\,III. The two
accounts reach the same hierarchy from opposite starting points: Memory Evolutive Systems take the
hierarchical category and the colimit as their apparatus, while the present account derives the
hierarchy from a physical substrate, the lattice and the confirmation fold of Volume 1. The
independent arrival at the same structure is a corroboration of it.

Closer to the result of \S\,V, the conserved surface is an invariant carried across events, and
that a persisting thing is such an invariant has a precedent in the identity-through-time
literature. Simons takes a continuant to be an invariant among occurrents under an equivalence
relation, the occurrents the more basic category and the continuant what holds across them
(Simons 2000). The surface here is an invariant of that kind, carried across the folds an
un-forming releases. The accounts differ at the ground: the invariant here is derived as the
compressible area-law surface (\S\,IV), not posited, so that the circularity pressed against bundle
accounts --- that a thing and its constituents must individuate each other (Lowe 1998) --- does not
arise here, the surface being fixed by the ledger independently of what it individuates.

Most directly, the question of whether the whole is a subject has been taken up in the
cosmopsychist program. Shani severs phenomenal constitution from phenomenal inclusion and argues
for a generative monism, on which the experiences of ordinary subjects are grounded in an
all-pervading cosmic consciousness without a mereological combination of parts into a whole
(Shani 2022). The present account runs toward the same whole from the opposite direction:
cosmopsychism grounds individual subjects downward from a cosmic subject, while this paper builds
upward and finds the top to be the one level that is not a subject in the closure sense, having no
exterior against which to close. The two roads fork at the subject of the whole, granted there and
withheld here, and the fork is located in the no-codomain result rather than in any prior
commitment about the whole.

\section*{Conclusion: A Boundary, Not a Mind}
\addcontentsline{toc}{section}{Conclusion: A Boundary, Not a Mind}

The confirmation tower has a top, and the operator that promotes a dissolved closure one level up is
defined everywhere on it but there. At the top the promotion has an argument and no recipient: there
is no level above to read the promoted symbol, and equally none to receive the two bits and the
volume residue the top's fold must shed. The empty codomain and the unreceived residue are one
absent recipient, read on the identity axis and the substrate axis of an un-forming. This is the
whole of what the top is, structurally --- not a larger subject and not a stronger one, but the
unique level at which the operation that defines every interior level runs out of the thing it
requires.

What follows from the absence is not a terminus. By the ledger that closes every interior fold, the
residue a fold sheds is the ground the next level is confirmed upon; applied at the top, where no
level waits to receive it, the residue constitutes that ground rather than finding it. The top
re-forms, by the same lossy operation that promotes a self at its death, carrying forward the
compressible surface --- the program identity, the shape of the whole --- and returning the interior
to the field as the ground of what follows. Nothing here is special to the top except the absence of
a recipient; the re-forming is the interior ledger applied where its recipient is not given but made.

And what the result does not touch, it does not touch. Whether the top closure is occupied ---
whether there is anything it is like to be the whole, in its standing or its re-forming --- is not
settled by where its promotion fails to go, nor by what its surface carries forward, and is not made
nearer by either. The ledger balances without that term: occupancy appears in neither the codomain,
nor the residue, nor the surface. The silence where a subject of the whole would be is not a gap in
the account but a boundary of it, and it is held.

What is settled here is narrower and firm. The top is a boundary condition on the promotion
operator: the one level at which promotion has an argument and no recipient, on both the identity
axis and the substrate axis at once. It does not dissolve, because its count is conserved; it does
not promote, because there is no level above; it re-forms, by the same lossy operation that promotes
a self at its death, carrying its surface forward and shedding its interior to the ground of what
follows. What is conserved across that re-forming is a surface --- a shape, an identity --- and the
question of whether that conserved thing has standing is the subject of the paper that follows, its
bridge imported there by hand. Here the result is the boundary and the re-forming: the top is the
level the operator cannot carry past, and the ledger carries it past anyway, by making the ground
its shedding had nowhere above to find.

\section*{Appendix A: The Un-forming Operator and Its Four Cases}
\addcontentsline{toc}{section}{Appendix A: The Un-forming Operator and Its Four Cases}

This appendix writes the operator by which a closure un-forms and shows its outcome set has exactly
four elements, each realized by a stated criterion. It inherits the boundary-reading of promotion
(Paper~7 \S\,II, Appendix~B.4; Paper~8 \S\,IV) and the conservation of bits (Axiom 3).

Let \ensuremath{C} be a level-\ensuremath{k} closure with constituent set \ensuremath{B(C)} and
boundary program \ensuremath{\pi(C)}, the symbol a level \ensuremath{k{+}1} reads (Paper~8 \S\,IV).
For a constituent \ensuremath{c \in B(C)}, write \ensuremath{\mathrm{cl}(c) = 1} if \ensuremath{c}
independently closes --- satisfies the level-\ensuremath{(k{-}1)} closure condition of Paper~7 \S\,II
on its own --- and \ensuremath{\mathrm{cl}(c) = 0} otherwise. For the program, write
\ensuremath{\mathrm{rd}(\pi) = 1} if some level \ensuremath{k{+}1} reads \ensuremath{\pi(C)} and
\ensuremath{\mathrm{rd}(\pi) = 0} if none does.

\subsection*{A.1 The two maps}

Define, on disjoint arguments by the orthogonality of Paper~7 Appendix~B.4,
\[
  \sigma(C) =
  \begin{cases}
    \text{recouple} & \text{if some } c \in B(C) \text{ with } \mathrm{cl}(c) = 1 \text{ enters a neighbouring closure,}\\
    \text{shed} & \text{otherwise,}
  \end{cases}
  \qquad
  \iota(C) =
  \begin{cases}
    \text{read} & \text{if } \mathrm{rd}(\pi) = 1,\\
    \text{lapse} & \text{if } \mathrm{rd}(\pi) = 0.
  \end{cases}
\]
The map \ensuremath{\sigma} reads the interior predicate \ensuremath{\mathrm{cl}}, a function of the
region alone; \ensuremath{\iota} reads \ensuremath{\mathrm{rd}}, a function of the surround alone. The
arguments are disjoint, so the two maps are independent. Both are total: by Axiom 3 a released bit
that does not enter a neighbour lies in \ensuremath{U_{\mathrm{sh}}}, so
\ensuremath{\sigma(C) \in \{\text{recouple}, \text{shed}\}} always; and \ensuremath{\mathrm{rd}} is
two-valued by excluded middle on the existence of a reader, so
\ensuremath{\iota(C) \in \{\text{read}, \text{lapse}\}} always.

\subsection*{A.2 The operator and its four cases}

The un-forming is the product map
\[
  U(C) = (\sigma(C), \iota(C)) \in \{\text{recouple}, \text{shed}\} \times \{\text{read}, \text{lapse}\},
  \qquad |\operatorname{im} U| = 2 \times 2 = 4.
\]
Each value is realized, fixed by a criterion on \ensuremath{\mathrm{cl}} and \ensuremath{\mathrm{rd}}:
\begin{itemize}
\item \ensuremath{(\text{recouple}, \text{read})} --- \emph{graft}: some constituent re-couples and a
  level reads \ensuremath{\pi(C)} (Paper~8 \S\,IV, the re-coupling outcome).
\item \ensuremath{(\text{shed}, \text{read})} --- \emph{promotion-in-death}: no constituent
  re-couples, \ensuremath{\pi(C)} is promoted; the death of Paper~8 (Conclusion).
\item \ensuremath{(\text{recouple}, \text{lapse})} --- \emph{decomposition}: constituents with
  \ensuremath{\mathrm{cl} = 1} re-couple, but \ensuremath{\mathrm{rd}(\pi) = 0} --- read as new
  structure, not as \ensuremath{\pi(C)}.
\item \ensuremath{(\text{shed}, \text{lapse})} --- \emph{total un-forming}:
  \ensuremath{\sigma(C) = \text{shed}} and \ensuremath{\mathrm{rd}(\pi) = 0}.
\end{itemize}
The partition is exhaustive because the product of two two-valued maps has exactly four values and
\ensuremath{U(C)} is single-valued.

\subsection*{A.3 Break-test}

The substrate map turns on the predicate \ensuremath{\mathrm{cl}}. A break would be a \ensuremath{C}
with \ensuremath{\sigma(C) = \text{shed}} yet possessing no constituent for which the shed is defined
--- a dissolution with neither recouple nor shed. This is impossible: if \ensuremath{\mathrm{cl}(c) = 0}
for every \ensuremath{c \in B(C)}, then no constituent independently closes, the binding-work
\ensuremath{C} performed on them is undone at un-forming, and that undone distinction is shed to
\ensuremath{U_{\mathrm{sh}}} by Axiom 3; so \ensuremath{\sigma(C) = \text{shed}} is well-defined
precisely there. Hence \ensuremath{\sigma} is total and no fifth value exists.

\section*{Appendix B: The Area/Volume Split and the Absent Recipient}
\addcontentsline{toc}{section}{Appendix B: The Area/Volume Split and the Absent Recipient}

This appendix computes what an un-forming carries and what it sheds, and derives the top's empty
codomain as the bit-side statement of \S\,II. It inherits the area law (Volume~2, Paper~1 \S\,2.2),
the per-fold shed (Paper~3 \S\,IV, Appendix~B), and the residue identification of Volume~1
(Paper~5 \S\,IV).

\subsection*{B.1 The two scalings}

A region \ensuremath{R} of linear extent \ensuremath{L} confirms on its frontier and acts through its
interior. By Volume~2 Paper~1 \S\,2.2 the now-surface frontier is forced to be a two-dimensional
agent manifold, so the accessible record --- the structure a level above can read --- scales as that
frontier's area,
\[
  I_{\mathrm{surface}}(L) \propto L^{d_s}, \qquad d_s = 2,
\]
the exponent fixed at \ensuremath{2} because it is the dimension of the agent manifold, not a fitted
constant (Volume~2, Paper~1 \S\,2.2; the coincidence with the Bekenstein exponent is each forced
from its own root, Paper~3 \S\,IV). The total fold activity ranges over the interior and scales as
the volume,
\[
  I_{\mathrm{total}}(L) \propto L^{d_v}, \qquad d_v = 3.
\]
The residue --- fold activity not surviving as accessible record --- is the excess of volume over
surface,
\[
  I_{\mathrm{residue}}(L) = I_{\mathrm{total}}(L) - I_{\mathrm{surface}}(L) \propto L^3 - L^2,
\]
positive and volume-dominant for \ensuremath{L > 1}, and provably uniform under
\ensuremath{U_{\mathrm{sh}}}'s own measure (Volume~1, Paper~5 \S\,IV).

\subsection*{B.2 What the split carries}

The surface term is the low-complexity part: a frontier of extent \ensuremath{L^2} admits a
description of order \ensuremath{L^2}, while the interior it bounds holds order \ensuremath{L^3} of
distinction, so
\[
  \frac{K(I_{\mathrm{surface}})}{K(I_{\mathrm{total}})} \sim \frac{L^2}{L^3} = \frac{1}{L} \to 0
  \quad (L \to \infty).
\]
The boundary is therefore what the level above reads --- the program identity \ensuremath{\pi(R)} of
Appendix~A is carried by \ensuremath{I_{\mathrm{surface}}} (Paper~7 \S\,II, Appendix~B.4: promotion
reads the boundary) --- and the interior \ensuremath{I_{\mathrm{residue}}} is what sheds to
\ensuremath{U_{\mathrm{sh}}}. This is the surface/interior split of \S\,V, computed: identity rides
the conserved \ensuremath{L^2} surface, lived interior returns to the field as the
\ensuremath{L^3 - L^2} residue.

\subsection*{B.3 The per-fold cost is local}

Each fold sheds a fixed two bits of distinction (Paper~3 \S\,IV): three binary loci span
\ensuremath{2^3 = 8} states, the confirmed outcome two, so
\[
  \text{shed per fold} = \log_2 8 - \log_2 2 = 3 - 1 = 2 \ \text{bits,}
\]
balanced against the fold's own three loci, retained bit, and two shed bits --- the fold's books
close at that fold. The quantity is the same at every fold (Paper~4 \S\,IV); no deficit compounds
across folds. The residue of \S\,B.1 is therefore not an accumulation of unpaid shed but the
per-level volume-minus-surface excess, shed and received at each level independently.

\subsection*{B.4 The absent recipient at the top}

At an interior level \ensuremath{k} the two shed bits and the volume residue are received by the
field \ensuremath{U_{\mathrm{sh}}} against which level \ensuremath{k{+}1} operates (Volume~1, Paper~5
\S\,IV); the fold closes. The top level has no level \ensuremath{k{+}1} (Paper~2 \S\,VI). Writing the
closure condition for a fold at level \ensuremath{k} as the existence of a recipient field for its
residue,
\[
  \text{fold closes at } k \iff \exists\, U_{\mathrm{sh}}^{(k+1)} \text{ receiving } I_{\mathrm{residue}}(L_k),
\]
the top fails the right-hand side: there is no \ensuremath{U_{\mathrm{sh}}^{(k+1)}}, so its residue
has no recipient. This is identically the empty codomain of \S\,II --- there \ensuremath{G(k)} has no
level \ensuremath{k{+}1} to receive the promoted symbol \ensuremath{\pi}; here the ledger has no level
\ensuremath{k{+}1} to receive the shed residue. The two are one absent recipient, read on the
identity axis and the substrate axis respectively. By Volume~1 Paper~5 \S\,IV the residue is itself
the lattice on which a level not previously operating is confirmed, so the top's residue, unreceived
by any existing level, constitutes the recipient field the closure condition requires --- which is
the re-forming of \S\,V, here the unique resolution of the closure condition at the top consistent
with the residue identification.

\subsection*{B.5 Break-test}

The no-codomain result of \S\,B.4 does not depend on the value of the exponent \ensuremath{d_s}.
Suppose \ensuremath{d_s} were not \ensuremath{2} but any \ensuremath{d_s < d_v}: the residue
\ensuremath{I_{\mathrm{residue}}(L) \propto L^{d_v} - L^{d_s}} remains positive and volume-dominant,
the surface remains the compressible part with
\ensuremath{K(I_{\mathrm{surface}})/K(I_{\mathrm{total}}) \to 0}, and the closure condition of \S\,B.4
is unchanged --- it asks only whether a recipient field exists, not how large the residue is. The
top fails that condition for the topological reason that there is no level above it (Paper~2
\S\,VI), independent of the exponent. The exponent \ensuremath{d_s = 2} fixes the \emph{magnitude}
of what is carried versus shed (\S\,B.1--B.2); it is the no-\emph{recipient} fact, not the
exponent, that makes the top special. A counterexample would therefore require not a different
exponent but a level above the top, which is excluded by the definition of the top. The result
stands on the absence of a recipient alone.

\section*{References}
\addcontentsline{toc}{section}{References}


\begin{reflist}
Ehresmann, A. C., \& Vanbremeersch, J.-P. (2007). \emph{Memory Evolutive Systems: Hierarchy, Emergence, Cognition}. Amsterdam: Elsevier.

Lambek, J. (1968). A fixpoint theorem for complete categories. \emph{Mathematische Zeitschrift}, 103(2), 151--161.

Lowe, E. J. (1998). \emph{The Possibility of Metaphysics: Substance, Identity, and Time}. Oxford: Clarendon Press.

Rutten, J. J. M. M. (2000). Universal coalgebra: A theory of systems. \emph{Theoretical Computer Science}, 249(1), 3--80.

Shani, I. (2022). Cosmopsychism, coherence, and world-affirming monism. \emph{The Monist}, 105(1), 6--24.

Simons, P. (2000). Identity through time and trope bundles. \emph{Topoi}, 19(2), 147--155.
\end{reflist}


\end{document}
